% Options for packages loaded elsewhere
\PassOptionsToPackage{unicode}{hyperref}
\PassOptionsToPackage{hyphens}{url}
\documentclass[
]{article}
\usepackage{xcolor}
\usepackage{amsmath,amssymb}
\setcounter{secnumdepth}{-\maxdimen} % remove section numbering
\usepackage{iftex}
\ifPDFTeX
  \usepackage[T1]{fontenc}
  \usepackage[utf8]{inputenc}
  \usepackage{textcomp} % provide euro and other symbols
\else % if luatex or xetex
  \usepackage{unicode-math} % this also loads fontspec
  \defaultfontfeatures{Scale=MatchLowercase}
  \defaultfontfeatures[\rmfamily]{Ligatures=TeX,Scale=1}
\fi
\usepackage{lmodern}
\ifPDFTeX\else
  % xetex/luatex font selection
\fi
% Use upquote if available, for straight quotes in verbatim environments
\IfFileExists{upquote.sty}{\usepackage{upquote}}{}
\IfFileExists{microtype.sty}{% use microtype if available
  \usepackage[]{microtype}
  \UseMicrotypeSet[protrusion]{basicmath} % disable protrusion for tt fonts
}{}
\makeatletter
\@ifundefined{KOMAClassName}{% if non-KOMA class
  \IfFileExists{parskip.sty}{%
    \usepackage{parskip}
  }{% else
    \setlength{\parindent}{0pt}
    \setlength{\parskip}{6pt plus 2pt minus 1pt}}
}{% if KOMA class
  \KOMAoptions{parskip=half}}
\makeatother
\usepackage{longtable,booktabs,array}
\usepackage{caption}
\captionsetup[table]{skip=6pt}
\newcounter{none} % for unnumbered tables
\usepackage{calc} % for calculating minipage widths
% Correct order of tables after \paragraph or \subparagraph
\usepackage{etoolbox}
\makeatletter
\patchcmd\longtable{\par}{\if@noskipsec\mbox{}\fi\par}{}{}
\makeatother
% Allow footnotes in longtable head/foot
\IfFileExists{footnotehyper.sty}{\usepackage{footnotehyper}}{\usepackage{footnote}}
\makesavenoteenv{longtable}
\setlength{\emergencystretch}{3em} % prevent overfull lines
\providecommand{\tightlist}{%
  \setlength{\itemsep}{0pt}\setlength{\parskip}{0pt}}
\usepackage{bookmark}
\IfFileExists{xurl.sty}{\usepackage{xurl}}{} % add URL line breaks if available
\urlstyle{same}
% fallback for those not using the hyperref driver hyperxmp:
\makeatletter
\@ifundefined{xmpquote}{\newcommand{\xmpquote}[1]{#1}}{}
\makeatother
\hypersetup{
  hidelinks,
  pdfcreator={LaTeX via pandoc}}

\author{}
\date{}

\begin{document}

\section{Gate E2 report --- run\_v1.2\_r1 (Paper IV
v1.2)}\label{gate-e2-report--run_v12_r1-paper-iv-v12}

Verdict: \textbf{INCONCLUSIVE}. Evidence directory:
\texttt{20\_EVIDENCE/E2/}.

\begin{center}\rule{0.5\linewidth}{0.5pt}\end{center}

\subsection{E2 --- Independent mathematics (first
pass)}\label{e2--independent-mathematics-first-pass}

Method: pen-and-paper rederivation by the lead auditor from the EN
manuscript (no author derivation notes, internal reports or editorial
reports were read); directed inspection of Lean sources only where the
prose is a summary (recorded). Exact numerical sub-steps are
cross-checked in E5 (tests T01--T24). Status legend: \textbf{R} =
rederived by the auditor; \textbf{R}* = rederived given an explicitly
stated intermediate estimate; \textbf{I} = imported with attribution
(not re-proved); \textbf{F} = only formally inspected (Lean), no human
derivation available/attempted; \textbf{X} = exposition gap (see A.2
table).

\subsection{Component ledger}\label{component-ledger}

{\def\LTcaptype{none} % do not increment counter
\begin{longtable}[]{@{}llll@{}}
\toprule\noalign{}
\# & Component (EN lines) & Status & Notes \\
\midrule\noalign{}
\endhead
\bottomrule\noalign{}
\endlastfoot
1 & Loss budget (1.3), (1.3a), Lemma 2.1, (2.4), Table 1 & R & pure
counting \\
2 & Finite LP duality / attained rational optimum (§2.2) & I/F & Paper
I; \texttt{PaperI.FiniteLPDuality} (not re-proved) \\
3 & (2.5) Paper II bound ⌊(2n+1)²/24⌋ = M(n) & R (floor identity), I
(Paper II theorem) & \\
4 & Lemma 3.2 bounded-rank nibble & I & Paper III, attributed \\
5 & Lemma 3.3 two quotas (l.1459-1483) & R & loads/codegrees of the
auxiliary extension, projection, both quota constants C=C0+1+k0+b(2+2k0)
rechecked \\
6 & Lemma 3.4 physical selector (l.1485-1504) & R & pair mass ≥
(t3−3)/2, load non-increase, codegree ≤ γ0/2+3/t3 (actual cross term
≤2/t3), gain 4m+b4, (3.6) \\
7 & (3.7) algebra; (3.9) & R & \\
8 & (3.10) charging (×3 per discarded edge; K4 with one bad edge keeps a
face; light K4 → faces at x/2) & R & discard count (3δ+1/k0+d)n² is
standard regularity bookkeeping, not derived (A2-2) \\
9 & (3.12)-(3.13) joint codegree & R* & given ψ\_H ≤ t² and fibre lower
bounds b\_H ≥ a3 t³, a4 t⁴ (A2-3) \\
10 & Per-edge load of the cleaned system (l.1615 "A single edge is
controlled in the same way") & X / F & \textbf{Following the stated
recipe gives load ≤ 1/a3 ≫ 1, not ≤ 1.} The Lean development instead
removes roots whose copy count deviates two-sidedly from the reference
value (\texttt{RC01DeviationCleanup.deviationBad}, Chebyshev
second-moment bound) and uses the resulting \emph{upper} spread bound
with budget (1+u)·A·pairVolume (\texttt{RC01CleanFiber} docstring). This
essential step is not stated in the prose. See A2-3. \\
11 & C.3 schedule (C.1), (C.2), Table C.1, (C.3) & F & numeric
domination with 2\^{}\{59516\}-scale quantities; definitions of L\_T,
M\_T, γ\_T, T\_T, D\_S, γ\_R are not in the manuscript (A2-4) \\
12 & Corollary 3.5 & R & uses M(n) ≥ n²/6 for n ≥ 6 \\
13 & Lemma 4.1 admissible mixed path & R & transport (4.4) (loads ≤2
before averaging, ≤1 after), existence of a non-twin step via Dirac +
rooted Dirac, lexicographic potential, chordality preservation \\
14 & Lemma 4.2, (4.5)-(4.6a), descent, near-threshold 4·10\^{}12 /
failure at 3.6·10\^{}12 & R & no circularity: Lemma 4.2 uses Thm 5.0,
whose proof does not use Prop 4.3 \\
15 & Prop 4.3 (4.7)-(4.8) & R & from \textbar ΔF4\textbar{} ≤ 6 per edit
(actually ≤4) and the first comparator branch \\
16 & (5.1a) three comparator branches & R (values of the branch maxima,
e.g. (4n+1)²/120), F (existence of the explicit fractional packings of
the middle/third branches) & \\
17 & Lemma 5.1 extraction & R & PEO edge bound (a−1)r−C(a,2); u ≤
√(2ε0)n; δ\_C ≤ 6.1ε0n²; a ≥ 0.33n; budget 1.6617e-6 n² \\
18 & Prop 5.2 regularization, (5.2)-(5.4), (5.7) & R & incl.
\textbar X\textbar{} ≤ a/16384, \textbar Y\textbar{} ≤ a/57344, exterior
cliques ≤ a/128+1, Y clique, R clique; d\_max: removed X-columns miss ≥
a/4 of ≤ 129a/64 exterior vertices ⇒ ≤ 113a/64 (the source of "113" is
not stated in the prose); all (5.4) ratios and (5.7) with margin \\
19 & Phase I: equitable Vizing colouring, (5.5), PEO orientation,
invalid incidences ≤ (w−1)A, cyclic averaging (5.6) & R & \\
20 & Phase II: (5.8)-(5.11) moments and averaging; compatibility with
phase I & R & (a+u)² ≤ Da+4ta+2tu \\
21 & (5.12)-(5.14) calibration; Table 2; Lemma 5.3 (5.15), (5.15a);
(5.16) & R & coefficients recomputed exactly (also E5 T03/T04) \\
22 & Corollary 5.4 & R* & given (4.8) \\
23 & Theorem B assembly; §6.1 weight argument; (6.2); §6.2 (6.3)-(6.4);
§6.3 Thm A with b=N² & R & tower value T(h+7) is F (item 11) \\
24 & Quadratic/linear coefficient optimality (6.5a), (6.5) & R & \\
25 & §6.4 (6.6) piece identity & R & \\
26 & Thm 6.1, (6.7a)-(6.7b), Cor 6.1a (48, 480), Cor 6.2 & R & 48 = best
from (6.7) (E5 T22); ≤10·d\_R bound (E5 T08) \\
27 & Thm C′: deletion step (6.15), δ′ ≥ 0, δ′ ≤ δ−εt+1/3, reinsertion ≤
t−1 edits paid by A\_s ≥ 2/ε\_s; window invariant (F.1 l.2141-2146);
repair bound B+4d\_E; (6.16) & R & invariant algebra also E5 T19 \\
28 & Thm C′ terminal accounts (6.14) / E.4 combination & R & coefficient
1+32s+800(3+32s)(s+1)² ≤ 25632(s+1)³ ≤ 38000(s+1)³; total ≤
40000(s+1)³ \\
29 & Cor 6.3 resize, (6.19) discrete parabola, (6.18) & R & parabola:
Q\_s−B = {[}(N−3k)(N−3k+1) − r\_N{]}/6, N=n+s, r\_N∈\{0,2\}; ≥ dist² (E5
T09) \\
30 & Prop 6.3a (6.20)-(6.21) & R & rsd ≤ s of G\_\{s,q,d\} (hosts: omit
\textbar D∩U\textbar; C, D vertices: host-clique argument), baseline
deficit (3d²−d)/2 via increments −1−3j, edge-count difference
d(4q−4s−d−1)/2, 4Δe − nd = d(5q−7s−2d−2) ≥ 0 \\
31 & Prop 6.4 (6.22)/(6.23) & R* (low-star branch: exact identity
(n−2c+1)c+C(c,2) = B\_\{n−s\}(c)+sc; (F.3)); I ({[}15, Lemma 3.1{]}
signed-star estimate (F.2), attributed); F (high-star envelope (F.4),
small orders n\textless8 in Lean) & independent finite check E5 T15 \\
32 & Thm 6.5 & R* & given SublinearEdit (from {[}22{]} via E34, I/F),
Lemma F.3a (R), (F.6) (R* --- my reconstruction gives 16δ+65t+n(s+u) and
δ+4t+2n(s+u)), RootScore minimization (R) \\
33 & Lemma F.3a, Cor F.3b & R & both induction cases; also E5 T12 \\
34 & Cor 6.6, (C.4) & F & tower comparisons (E35) not rederived \\
35 & Appendix E.1 constructor: (E.2), (E.2a), (E.2c), (E.2e) & R &
(E.2e) increments ⌊(N+s+1)/3⌋ ≥ ⌈(c+b+s)/3⌉ \\
36 & (E.1) capacity conditions, list hosting of all edges of G{[}S{]}
(Galvin, dyadic) & F & only described \\
37 & (E.2b) three-case budget & R* & cases 1--3 arithmetic checked given
(E.2d) and the normalization estimates \\
38 & (E.2d) joint defect bound & X & sketch only (A2-6) \\
39 & Normalization estimates (t\_W, L, \textbar 𝓛\textbar,
\textbar c−n/3\textbar, σ, τ, m) & X / F & asserted, not derived
(A2-5) \\
40 & E.2 min-degree reduction removing K\_s; (E.3) & R & (E.3) holds for
all N ≥ 1 \\
41 & E.3 lower witness, (E.5) & R & \\
42 & E.4 block by block: StrictParams (κ, q, capacity conditions) & R*
(given normalization numbers) & κ ≥ 2n/3 − 0.021ρ; light/heavy host
degree bounds ≤ κ; 3⌈c/2⌉+2τ+4q ≤ b+2 \\
43 & E.4 (E.7) θ\_s-saving from cyclic averaging & R* & each H-edge has
≥ c−2σ valid hosts; z ≥ dm/κ ≥ θ\_s m \\
44 & E.4 (E.8)-(E.9) & R & exact substitution of R\_* and F; D\_0 ≥ 0
from (E.2e) \\
45 & E.4 StrictBudget (E.10) & R* & depends on (E.2d)+normalization
(items 38--39) \\
46 & E.4 ResidualExcess (E.11) & R & needs 2(L+t\_W) ≤ (c+b+s)/3 \\
47 & E.4 missing-pair restriction (no\_core\_pairs) & R* & given
\textbar U\textbar{} ≥ ρ/10⁴+s+2 from normalization \\
48 & E.4 CoreCliqueAlternative (E.12), (E.13) & R & induction on s;
(4s+1)s missing-edge bound (4s−1 suffices); c ≥ 20(s+1)² \\
49 & E.4 TemplateDistance / RootPartition (E.14) & R & two successive
comparisons \\
50 & F.5 case structure A/B1/B2 & R (assembly logic, exhaustiveness,
(F.12) with δ11 ≤ ε/2, Kη = ε/2, prefix-extension counting) ; I/F
(adapted set-colouring theorem thm2, claim5\_count, gluing lemma ---
from {[}22{]} with modifications) & matches Lean \texttt{lemma11\_V} \\
51 & (F.7)-(F.10) auxiliary lemmas & F (+ E5 finite tests for F.7, F.9)
& \\
52 & Prop D.1, Table 8, (D.1)-(D.4) & R (counts, residue coverage), I
(Bose/Skolem STS existence) & \\
53 & Prop D.2, D.3 (D.6)-(D.7), D.3 (D.8) & R & \\
54 & G.1 annex (G.1) & F (statement only) & \\
55 & Prop G.1 (G.2)-(G.5) & R & Behrend via Mathlib is I \\
56 & §8 (8.1) & R & \\
\end{longtable}
}

No mathematical error was found in any rederived step.

\subsection{Appendix A.2 exposition verdicts (mandate
§5.2)}\label{appendix-a2-exposition-verdicts-mandate-52}

{\def\LTcaptype{none} % do not increment counter
\begin{longtable}[]{@{}lllllll@{}}
\toprule\noalign{}
Item & Manuscript location & Literal Lean declarations (located) &
Independent derivation attempted & Evidence & Exposition verdict &
Explanation \\
\midrule\noalign{}
\endhead
\bottomrule\noalign{}
\endlastfoot
A2-1 & §5.1 (5.3), (5.4), (5.7) l.484-514 & \texttt{Regularization},
\texttt{RegularizationBounds}, \texttt{RegularizedRootGoal},
\texttt{RootRegularizationBridge} & Yes, complete (component 18) & this
record; E5 T06 & \textbf{ACCEPTABLE\_SUMMARY} & Every bound follows from
(5.1), (5.3) and the displayed degree/width counts by rational
comparison with wide margins; the one non-obvious input (removed
X-columns miss ≥ a/4 exterior vertices, giving 113a/64) should be one
sentence. \\
A2-2 & C.2 (3.10), (3.8) l.1572-1581, l.1538-1544 &
\texttt{E17.Cleanup}, \texttt{E17.FarBudget},
\texttt{RC01RootwiseCleanupBudget}, \texttt{RC01RootwiseRetention} &
Partial: charging argument rederived; discard count reconstructed from
standard regularity bookkeeping (exceptional class, irregular pairs,
intra-class, sparse pairs) & component 8 & \textbf{ACCEPTABLE\_SUMMARY}
for the discard counts (3.10); retention (3.8) is covered under A2-3 & A
referee familiar with regularity can reconstruct (3.10); the conventions
(δ/8 in C.3, class-size t) are stated. \\
A2-3 & §3.2, C.2 fibre bounds b\_H ≥ a3t³, a4t⁴; (3.8); per-edge
feasibility l.1606-1615 & \texttt{RC01CleanFiber},
\texttt{RC01DeviationCleanup}, \texttt{RC01RegularVolume},
\texttt{RC01PatternMassScale}, \texttt{RC01CleanedGate},
\texttt{RC01CleanedMixedCodegree} & Yes; the codegree bound (3.13)
follows; the per-edge load bound does \textbf{not} follow from the text
& component 10 & \textbf{REQUIRES\_EXPANSION} & The feasibility (load ≤
1) of the cleaned system needs a two-sided deviation cleanup of roots
with a second-moment (Chebyshev) bound and an upper spread bound; the
prose instead says "a single edge is controlled in the same way", which
with the stated bounds yields only 1/a3. The retention (3.8) likewise
depends on that cleanup and the counting lemma. This is an essential
step of Theorem 3.1 absent from the human-readable proof (the formal
proof contains it). \\
A2-4 & C.3 Table C.1, (C.1)-(C.3) & \texttt{E19.ScheduleBounds},
\texttt{E19.ChainBounds}, \texttt{E19.GateBounds}, \texttt{E18.Numeric},
\texttt{E17.ExplicitFarAssembly} & No (definitions of the schedule
functions are not given in the text) & --- &
\textbf{REQUIRES\_EXPANSION} (minor impact) & Affects only the explicit
tower values T(h+7), T(h+8) and (C.4)/(6.28), not the existence of
thresholds. A human referee cannot check Table C.1 from the manuscript;
either give the schedule definitions or label these values as formally
verified numerics only. \\
A2-5 & E.1 normalization l.1853-1863; E.4 l.1993-2003 &
\texttt{A4S1.IndepAll*} (AllInput.params\_strict, normalization
producers), E33 localization & Arithmetic downstream checked (items
42--47); the normalization estimates themselves could not be derived
from the text & components 39, 42--47 & \textbf{REQUIRES\_EXPANSION} &
These estimates (t\_W ≤ 4ρ/10\^{}27, L ≤ ρ/10\^{}13+1,
\textbar 𝓛\textbar−n/3 ≤ c+4λ/10\^{}31, \textbar c−n/3\textbar, σ, τ, m
≤ 2ρ²/10\^{}31) carry the whole fixed-defect constructor, i.e. the proof
of Theorem C for s\textgreater0 (the paper\textquotesingle s order-four
strengthening relative to {[}15{]}) and Theorem C′. They are asserted,
with the localization input itself only named. \\
A2-6 & E.1 (E.2b), (E.2d), three-case budget l.1837-1929 &
\texttt{IndepAllPeel}, \texttt{IndepAllObstr}, \texttt{IndepAllBudget},
\texttt{IndepAllMain} & Cases 1--3 checked given (E.2d); (E.2d) itself
only sketched & components 37--38 & \textbf{REQUIRES\_EXPANSION} &
(E.2d) is the key place where rooted defect is used jointly on up to
2s+2 exceptions; the text gives an outline ("the rooted condition,
applied to the induced subgraphs determined by these pairs \ldots{}
bounds the joint contribution") that a referee cannot turn into a proof
without substantial reconstruction. \\
A2-7 & D.1 Table 8 l.1721-1744 & \texttt{HalvingFrames},
\texttt{ThreeRegime.CompleteStateAllOrders}, \texttt{exists\_packing},
\texttt{complete\_state\_closes} & Yes: residue coverage and gain counts
(D.1)-(D.4) rederived; STS existence classical & component 52 &
\textbf{ACCEPTABLE\_SUMMARY} & Classical Bose/Skolem Steiner triple
systems with a standard parallel-class K4 augmentation; Proposition D.1
is a supplement not used by Theorems A--C. \\
\end{longtable}
}

\subsection{§5.1 obligations (v1.2-specific) --- first-pass
results}\label{51-obligations-v12-specific--first-pass-results}

\begin{itemize}
\tightlist
\item
  \textbf{E.4 block by block}: StrictParams R*; StrictBudget R*
  (conditional on A2-5/A2-6 items); ResidualExcess R;
  CoreCliqueAlternative R; LargeCliqueCore / no\_core\_pairs R*;
  TemplateDistance R; RootPartition R. Same witnesses: all accounts
  refer to one physical partition Q and one normalized triple (S,H,W);
  (E.9) uses the constructed Q, not an optimizer. Fixed-s thresholds and
  windows: γ\_s, N\_s\^{}stab match (F.1)/(6.11) term by term (E3).
\item
  \textbf{F.3a}: \texttt{exists\_clique\_additive\_defect} ---
  hypotheses rsd ≤ s on G, arbitrary D ⊆ V, output clique C ⊆ D of G
  with (\textbar D\textbar−\textbar C\textbar−s)\_+² ≤ ordNE(G,D):
  genuine clique recovery (not an induced-subgraph or template
  assertion). R.
\item
  \textbf{F.5}: l11\_caseA / B1 / B2 conditions, exhaustiveness and
  common conclusion checked (E3 static; component 50). Adaptation vs
  {[}22{]}: sequences with replacement, K = nodes, enlarged pinned
  class, constants m11, B, δ11; accurate per E7 lead verification.
\item
  \textbf{Theorem C′ quantifier}: one root (C,D,H) chosen before every Q
  and τ (Lean ∃CDH \ldots{} ∀Q τ). Linear bound is to an actual-root
  template of arbitrary core size; optimal-size comparison carries
  n√((1+4A\_s)δ) (Cor 6.3) and Prop 6.3a shows the square-root scale is
  necessary along an explicit family for each fixed s. No growing s is
  substituted.
\item
  \textbf{Uniform/sublinear}: Prop 6.4 threshold independent of s; Thm
  6.5 via SublinearEdit/{[}22{]} and RootScore; E35 uniform range ∀s ≤ S
  at n ≥ T(P(S)). Checked statically.
\item
  \textbf{Selected Theorem C chain}:
  \texttt{DefectExplicitPublication.rooted\_defect\_eventual} :=
  ⟨E33.Fexp E34.NeditE s, E34.theoremC\_fully\_explicit\_final⟩; E34
  chain uses \texttt{lemma11\_V} (F.5). Alon--Shapira exclusion on the
  six public cones is a dynamic check (E4).
\item
  \textbf{B7 / tower}: (3.11) link and (C.1)-(C.3) are formally stated;
  tower propagation not humanly rederivable from the text (A2-4).
  Inverse threshold: \texttt{sMaxT} requires n ≥ T(P(0)) (stated).
\item
  \textbf{gainOf conventions}: \texttt{LossBudget.budget\_iff} and
  \texttt{SplitMixedGap.mixed\_gap\_zero} use
  \texttt{FarRounding.gainOf} (C(\textbar K\textbar,2)∸1);
  partition-stability arguments use edge/defect accounting. Matches
  §7.2.
\end{itemize}

\subsection{Gate E2 first-pass
verdict}\label{gate-e2-first-pass-verdict}

\textbf{INCONCLUSIVE} --- no mathematical error found; all headline
assemblies rederived; but four A.2 items (A2-3, A2-4, A2-5, A2-6) are
essential or quantitative steps whose human-readable derivation is
absent or only sketched, and one of them (A2-3) is stated in a way that
does not by itself give the claimed feasibility. Their correctness
therefore rests on the Lean development (to be confirmed in E4) and not
on the manuscript. These are exposition findings, not counterexamples.

\end{document}
